\section{Khung pháp lý về bảo vệ dữ liệu bệnh nhân tại Việt Nam}
\label{sec:legal}

Hình ảnh y tế là dữ liệu cá nhân thuộc nhóm nhạy cảm nhất, và tại Việt Nam việc
xử lý loại dữ liệu này chịu sự điều chỉnh của một khung pháp lý nhiều tầng gồm
pháp luật chuyên ngành y tế, pháp luật chung về bảo vệ dữ liệu cá nhân và pháp
luật về an ninh mạng. Các ràng buộc này mang tính kiến trúc: chúng quyết định
dữ liệu được lưu trú ở đâu, ai được phép truy cập và những gì bắt buộc phải ghi
nhận, và không thể bổ sung thêm về sau vào một hệ thống vốn được thiết kế mà
không tính đến chúng.

\subsection{Pháp luật chuyên ngành y tế}

\textbf{Luật Khám bệnh, chữa bệnh} (số 15/2023/QH15, có hiệu lực từ ngày 01
tháng 01 năm 2024) \cite{lawmet2023} xác lập quyền được bảo mật thông tin của người bệnh như một
quyền luật định. Thông tin về tình trạng sức khoẻ và đời tư của người bệnh được
ghi trong hồ sơ bệnh án phải được giữ bí mật và chỉ được tiết lộ khi có sự đồng
ý của người bệnh hoặc trong các trường hợp pháp luật cho phép. Cũng chính luật
này điều chỉnh bản thân hồ sơ bệnh án, quy định thời hạn lưu trữ và xác định
những chủ thể được phép tiếp cận, đồng thời khẳng định chẩn đoán là hoạt động
chỉ dành cho người hành nghề có giấy phép. Điểm cuối cùng này là cơ sở pháp lý
cho định hướng hỗ trợ được giữ nhất quán trong toàn bộ báo cáo: nền tảng có thể
cung cấp thông tin phục vụ chẩn đoán nhưng không được đưa ra chẩn đoán.

\textbf{Thông tư 13/2025/TT-BYT} hướng dẫn triển khai hồ sơ bệnh án điện tử
(ban hành ngày 06/6/2025, có hiệu lực từ ngày 21/7/2025) thay thế Thông tư
46/2018/TT-BYT \cite{tt13_2025,mohemr}. Thông tư đặt lộ trình bắt buộc: các bệnh
viện phải triển khai chậm nhất ngày 30/9/2025, các cơ sở khám bệnh, chữa bệnh khác
chậm nhất ngày 31/12/2026. Cơ sở phải có hạ tầng gồm máy chủ, thiết bị lưu trữ,
giải pháp an toàn thông tin và phần mềm quản lý hồ sơ; hệ thống phải đáp ứng tiêu
chuẩn kỹ thuật về an toàn thông tin và có khả năng phục hồi dữ liệu, hỗ trợ chữ ký
số hoặc phương thức xác thực hợp pháp khác, và kết nối với cơ sở dữ liệu quốc gia
về dân cư. Hai yêu cầu có hệ quả kiến trúc trực tiếp. Yêu cầu truy xuất và kiểm tra
hồ sơ về sau khiến một nhật ký truy vết chỉ ghi thêm (append-only), ghi nhận mọi
truy cập, trở thành yêu cầu chức năng chứ không còn là lựa chọn thiết kế. Yêu cầu
khả năng phục hồi dữ liệu biến độ bền của kho ảnh thành nghĩa vụ pháp lý, không
chỉ là chất lượng kỹ thuật.

\subsection{Bảo vệ dữ liệu cá nhân}

Khung bảo vệ dữ liệu cá nhân của Việt Nam đã thay đổi căn bản trong thời gian
thực hiện đề tài. \textbf{Nghị định 13/2023/NĐ-CP}, có hiệu lực từ ngày 01/7/2023,
là văn bản toàn diện đầu tiên \cite{decree13}. Kể từ ngày 01/01/2026, vai trò này
thuộc về \textbf{Luật Bảo vệ dữ liệu cá nhân số 91/2025/QH15}, được Quốc hội thông
qua ngày 26/6/2025, cùng \textbf{Nghị định 356/2025/NĐ-CP} quy định chi tiết thi
hành, thay thế Nghị định 13 \cite{pdpl2025,decree356}. Các hồ sơ đánh giá tác động
đã lập theo Nghị định 13 vẫn được công nhận, nhưng mọi cập nhật về sau phải tuân
theo luật mới. Phần lớn nghĩa vụ được kế thừa từ Nghị định 13, nay được nâng lên
tầm luật và đi kèm chế tài nặng hơn. Một số quy định tác động trực tiếp đến thiết
kế hệ thống.

Luật phân biệt \emph{dữ liệu cá nhân cơ bản} với \emph{dữ liệu cá nhân nhạy cảm}
và xếp thông tin về tình trạng sức khoẻ vào nhóm nhạy cảm, kéo theo các nghĩa vụ
tăng cường. Sự đồng ý phải tự nguyện, rõ ràng và cụ thể cho từng mục đích; sự im
lặng không được coi là đồng ý, và chủ thể dữ liệu có quyền rút lại sự đồng ý bất
kỳ lúc nào. Trên thực tế, bên kiểm soát dữ liệu phải lưu giữ được bằng chứng về sự
đồng ý cùng trạng thái hiện hành của nó. Chủ thể dữ liệu được trao các quyền truy
cập, chỉnh sửa, xoá và phản đối, với thời hạn phản hồi được quy định cụ thể. Bên
kiểm soát dữ liệu phải lập hồ sơ \textbf{đánh giá tác động xử lý dữ liệu cá nhân}
và gửi cơ quan chuyên trách bảo vệ dữ liệu cá nhân thuộc \textbf{Bộ Công an} trong
thời hạn 60 ngày (Điều 21).

Đối với \textbf{chuyển dữ liệu cá nhân xuyên biên giới} (Điều 20), bên chuyển dữ
liệu phải lập hồ sơ đánh giá tác động chuyển dữ liệu và gửi cơ quan chuyên trách
trong thời hạn 60 ngày. Vi phạm quy định về chuyển dữ liệu xuyên biên giới có thể
bị phạt tới 5\% doanh thu của năm trước liền kề, các vi phạm khác tới 3 tỷ đồng.
Đối với dự án này, hệ quả là trực tiếp và nặng nề: một kiến trúc lưu trữ hình ảnh
trên dịch vụ đám mây đặt ở nước ngoài sẽ phát sinh gánh nặng tuân thủ và rủi ro chế
tài mà một triển khai tại chỗ khép kín hoàn toàn tránh được. Chính cân nhắc này, hơn
bất kỳ yếu tố kỹ thuật nào, quyết định mô hình ghi cục bộ trước và đặt mọi kho ngoài
cơ sở trong lãnh thổ Việt Nam được đề xuất trong Mục~\ref{sec:architecture}.

Song song, \textbf{Luật Dữ liệu số 60/2024/QH15}, có hiệu lực từ ngày 01/7/2025,
phân loại dữ liệu thành dữ liệu cốt lõi, dữ liệu quan trọng và dữ liệu khác, và đặt
ra thủ tục riêng đối với việc chuyển và xử lý dữ liệu cốt lõi, dữ liệu quan trọng ra
nước ngoài \cite{datalaw2024}. Một kho ảnh tập trung quy mô vùng có thể rơi vào các
nhóm này khi danh mục được ban hành. Nguyên tắc thiết kế rút ra là xem các yêu cầu
hiện hành như mức sàn chứ không phải mức trần.

\subsection{An ninh mạng và nội địa hoá dữ liệu}

\textbf{Luật An ninh mạng} số 24/2018/QH14 \cite{cybersecuritylaw} và
\textbf{Nghị định 53/2022/NĐ-CP} hướng dẫn thi hành \cite{decree53} đặt ra nghĩa vụ
nội địa hoá dữ liệu đối với một số nhóm nhà cung cấp dịch vụ, yêu cầu các loại dữ
liệu được quy định về người dùng Việt Nam phải được lưu trữ trong lãnh thổ Việt Nam.
Từ ngày 01/7/2026, luật này được thay thế bởi \textbf{Luật An ninh mạng số
116/2025/QH15} \cite{cybersecuritylaw2025}, mở rộng phạm vi bảo vệ tới hệ thống
thông tin và dữ liệu, và yêu cầu chủ quản hệ thống thực hiện quản lý rủi ro, kiểm
tra an ninh, ứng cứu sự cố và các giải pháp bảo vệ dữ liệu, đặc biệt là dữ liệu cá
nhân. Các văn bản hướng dẫn luật cũ, trong đó có Nghị định 53/2022/NĐ-CP, cần được
rà soát lại khi văn bản hướng dẫn luật mới được ban hành; về phía thiết kế, giả định
an toàn là dữ liệu bệnh nhân được lưu trữ trong lãnh thổ Việt Nam. Hệ thống
thông tin bệnh viện còn có thể thuộc phạm vi điều chỉnh của các quy định về hệ
thống thông tin theo cấp độ, theo đó các yêu cầu bảo đảm an toàn được áp dụng
tương ứng với mức độ quan trọng được xác định của hệ thống. Một nền tảng can
thiệp vào quy trình chẩn đoán cấp cứu cần được mặc định xem là thuộc cấp độ
không thấp, và phải được thiết kế phù hợp với giả định đó.

\subsection{Các yêu cầu thiết kế phát sinh}

Phân tích pháp lý dẫn tới các yêu cầu được chuyển trực tiếp sang
Mục~\ref{sec:architecture}:

\begin{enumerate}[label=L\arabic*.,leftmargin=1.2cm]
    \item \textbf{Nơi lưu trú dữ liệu.} Dữ liệu hình ảnh định danh được mặc
          định giữ trong phạm vi bệnh viện, hoặc trong lãnh thổ Việt Nam. Mọi
          đường truyền dữ liệu ra nước ngoài đều đòi hỏi đánh giá tác động
          chuyển dữ liệu và được xem là trường hợp ngoại lệ.
    \item \textbf{Tối thiểu hoá dữ liệu.} Chỉ những dữ liệu thực sự cần thiết
          cho từng mục đích mới được truyền ra khỏi kho lưu trữ gốc; các thuộc
          tính định danh được loại bỏ ngay tại ranh giới theo Mục~\ref{sec:deid}.
    \item \textbf{Kiểm soát truy cập.} Xác thực là bắt buộc và theo vai trò;
          cấu hình cho phép truy cập ẩn danh tuy chấp nhận được trong bản mẫu
          phát triển nhưng bị nghiêm cấm khi triển khai thực tế.
    \item \textbf{Truy vết.} Mọi hành vi truy cập và tiết lộ dữ liệu bệnh nhân
          đều được ghi vào nhật ký chỉ ghi thêm, với mức chi tiết đủ để tái lập
          ai đã truy cập dữ liệu nào và vào thời điểm nào.
    \item \textbf{Mã hoá.} Dữ liệu được mã hoá khi truyền, điều mà DICOM cổ
          điển đòi hỏi phải dùng TLS hoặc đường hầm, và mã hoá khi lưu trữ.
    \item \textbf{Thời hạn lưu trữ.} Thời hạn lưu trữ có thể cấu hình theo yêu
          cầu áp dụng, kèm cơ chế xoá có căn cứ khi hết hạn.
    \item \textbf{Sự đồng ý và giới hạn mục đích.} Dữ liệu thu thập phục vụ
          khám chữa bệnh không được tái sử dụng cho mục đích nghiên cứu nếu
          không có căn cứ pháp lý riêng, và hệ thống phân tách hai mục đích này
          bằng cấu trúc chứ không chỉ bằng chính sách.
\item \textbf{Khả năng phục hồi.} Hình ảnh, như một phần của hồ sơ bệnh án
      điện tử, phải khôi phục được sau sự cố theo Thông tư 13/2025/TT-BYT;
      yêu cầu này được hiện thực bằng các cơ chế độ bền ở
      Mục~\ref{sec:durability}.
\end{enumerate}

\medskip